\pdfvariable trailerid {[<C3182026090300000000000000000000><C3182026090300000000000000000000>]}
\documentclass[10pt]{article}
\usepackage[margin=0.72in]{geometry}
\usepackage{amsmath,amssymb,amsthm,booktabs,microtype,hyperref}
\hypersetup{colorlinks=true,linkcolor=blue,urlcolor=blue}
\providecommand{\CRevisionRound}{2}
\newtheorem{theorem}{Theorem}
\newtheorem{corollary}[theorem]{Corollary}
\newcommand{\sinc}{\operatorname{sinc}}
\newcommand{\spec}{\operatorname{spec}}
\newcommand{\iu}{\mathrm{i}}
\setlength{\emergencystretch}{3em}
\title{An Exact Finite-Size Bulk--Edge and Quench Atlas for the SSH Chain}
\author{Route-A source-local certificate HCS-C318}
\date{3 September 2026}

\begin{document}
\maketitle
\begin{abstract}
For the balanced open Su--Schrieffer--Heeger chain of $M$ cells, we derive
the full characteristic polynomial and prove that a strictly hyperbolic
edge pair exists if and only if $w/v>(M+1)/M$.  We give its exact splitting,
two-ended eigenvectors, decay bound, critical linear taper, and fixed-ratio
asymptotic.
\ifnum\CRevisionRound>0
We then separate this finite threshold from the periodic bulk transition,
resolve finite-ring parity and every hopping face, and derive an entire
matrix-sinc propagator valid at singular couplings.
\fi
\ifnum\CRevisionRound>1
An exact mode-quench corollary retains the finite momentum-grid caveat;
independent symbolic, numerical, replay, and hostile tests define the claim
boundary.
\fi
\end{abstract}

\section{Finite hyperbolic edge threshold}
Let $M\geq2$ and $v,w\geq0$.  In the ordered basis
$(A_1,\ldots,A_M,B_1,\ldots,B_M)$, the open Hamiltonian is
\begin{equation}\label{eq:block}
 H_O=\begin{pmatrix}0&T\\T^*&0\end{pmatrix},\qquad
 T=vI+wS_-,\qquad (S_-)_{j,j-1}=1.
\end{equation}
Thus $\Gamma=\operatorname{diag}(I,-I)$ obeys
$\Gamma H_O\Gamma=-H_O$.

\begin{theorem}[Characteristic and sharp finite edge theorem]\label{thm:edge}
Put $q_M(y)=\det(yI-TT^*)$.  Then
\begin{equation}\label{eq:recurrence}
 q_0=1,\quad q_1=y-v^2,\quad
 q_m=(y-v^2-w^2)q_{m-1}-v^2w^2q_{m-2}.
\end{equation}
For $v,w>0$, set $x=(E^2-v^2-w^2)/(2vw)$.  With $U_m$ the Chebyshev
polynomial of the second kind,
\begin{equation}\label{eq:characteristic}
 \det(EI-H_O)=(vw)^M
 \left[U_M(x)+\frac{w}{v}U_{M-1}(x)\right].
\end{equation}
There is exactly one secular root $x=-\cosh\kappa<-1$, and hence exactly
one eigenvalue pair $\pm E_{\rm e}$ outside the trigonometric sector, if and
only if
\begin{equation}\label{eq:threshold}
 \frac{w}{v}>\frac{M+1}{M}.
\end{equation}
For the unique $\kappa>0$ determined by
\begin{equation}\label{eq:kappa}
 \frac{w}{v}=\frac{\sinh((M+1)\kappa)}{\sinh(M\kappa)},
\end{equation}
the positive energy and one unnormalised positive-energy vector are
\begin{align}
 E_{\rm e}
 &=v\frac{\sinh\kappa}{\sinh(M\kappa)}
 =w\frac{\sinh\kappa}{\sinh((M+1)\kappa)},\label{eq:energy}\\
 a_j&=(-1)^{j-1}\sinh((M+1-j)\kappa),\qquad
 b_j=(-1)^{j-1}\sinh(j\kappa).\label{eq:vectors}
\end{align}
Thus $(a,b)$ has energy $E_{\rm e}$ and $(a,-b)$ has energy $-E_{\rm e}$.
The profiles decay strictly inward:
\begin{equation}\label{eq:decay}
 \frac{|a_{j+1}|}{|a_j|}<e^{-\kappa},\qquad
 \frac{|b_j|}{|b_{j+1}|}<e^{-\kappa}\quad(1\leq j<M).
\end{equation}
At equality in~\eqref{eq:threshold}, $x=-1$ and $E=v/M$.  After common
division of $a,b$ by $\kappa$ (equivalently, after overall normalisation),
the $\kappa\downarrow0$ profiles are
$a_j=(-1)^{j-1}(M+1-j)$ and $b_j=(-1)^{j-1}j$.
\end{theorem}

\begin{proof}
The leading diagonal entry of $TT^*$ is $v^2$; every subsequent one is
$v^2+w^2$, and the off-diagonal entries are $vw$.  Its continuant is
\eqref{eq:recurrence}.  Normalising this recurrence gives
\eqref{eq:characteristic}; in particular, the factor $(vw)^M$ cannot be
dropped.  The $M$ roots are real and simple because $TT^*$ is an irreducible
Jacobi matrix.  Also $\det T=v^M$, so $E=0$ is impossible when $v>0$.

For $x=-\cosh\kappa$, the identity
$U_m(-\cosh\kappa)=(-1)^m\sinh((m+1)\kappa)/\sinh\kappa$ reduces the secular
equation to~\eqref{eq:kappa}.  Its logarithmic derivative is
\[
 (M+1)\coth((M+1)\kappa)-M\coth(M\kappa)>0,
\]
because $s\mapsto s\coth(s\kappa)$ is strictly increasing.  The ratio
therefore rises from $(M+1)/M$ to infinity.  No root lies above $x=1$,
where both Chebyshev terms are positive, so this proves the iff statement
and uniqueness.  Direct hyperbolic addition verifies $Tb=E_{\rm e}a$ and
$T^*a=E_{\rm e}b$, including both expressions in~\eqref{eq:energy}.
Finally,
\[
 \frac{\sinh(n\kappa)}{\sinh((n+1)\kappa)}
 =e^{-\kappa}\frac{1-e^{-2n\kappa}}
 {1-e^{-2(n+1)\kappa}}<e^{-\kappa},
\]
which proves~\eqref{eq:decay}.  Taking $\kappa\downarrow0$ gives the
threshold energy; applying the common $1/\kappa$ rescaling before that limit
gives the nonzero tapers.
\end{proof}

\begin{corollary}[Finite splitting]\label{cor:splitting}
For fixed $r=w/v>1$ and all sufficiently large $M$, the hyperbolic root
exists, $\kappa_M\to\log r$, and
\begin{equation}\label{eq:asymptotic}
 E_{\rm e}\sim w(1-r^{-2})r^{-M}.
\end{equation}
Consequently the finite pair is never exactly zero for $v>0$, although its
splitting is exponentially small.
\end{corollary}
\begin{proof}
Writing $z=e^{-\kappa}$ turns~\eqref{eq:kappa} into
$r=z^{-1}(1-z^{2M+2})/(1-z^{2M})$, hence $z\to r^{-1}$.  Equation
\eqref{eq:energy} becomes
$E_{\rm e}=v z^{M-1}(1-z^2)/(1-z^{2M})$, which gives
\eqref{eq:asymptotic}.
\end{proof}

\ifnum\CRevisionRound>0
\section{Bulk winding, boundary faces, and unitary propagation}
Close the chain and use momenta $k_n=2\pi n/M$.  Our orientation convention is
\begin{equation}\label{eq:bloch}
 q(k)=v+we^{\iu k},\qquad
 h(k)=\begin{pmatrix}0&\overline{q(k)}\\q(k)&0\end{pmatrix}.
\end{equation}

\begin{theorem}[Periodic bulk and finite parity]\label{thm:bulk}
The Bloch energies are
\begin{equation}\label{eq:dispersion}
 E_\pm(k)=\pm\sqrt{v^2+w^2+2vw\cos k}.
\end{equation}
For the continuum symbol and the sampled $M$-cell ring, respectively, let
$\Delta_\infty$ and $\Delta_M$ denote the minimum positive-band distance
to zero.  Then
\begin{equation}\label{eq:gaps}
 \Delta_\infty=|v-w|,\qquad
 \Delta_M=
 \begin{cases}
 |v-w|,&M\ \text{even},\\
 \sqrt{v^2+w^2-2vw\cos(\pi/M)},&M\ \text{odd}.
 \end{cases}
\end{equation}
The corresponding central band gaps are $2\Delta_\infty$ and $2\Delta_M$.
Whenever the loop avoids zero, its counterclockwise winding is
$1$ for $w>v$ and $0$ for $v>w$; at $v=w>0$ it is undefined.  At this
critical line a finite ring has a two-dimensional zero sector if and only if
$M$ is even.  For odd $M$ it instead has
$\Delta_M=2v\sin(\pi/(2M))>0$.
\end{theorem}
\begin{proof}
Squaring~\eqref{eq:bloch} gives $|q(k)|^2I$, proving the dispersion.
Over the continuum, $\cos k$ is minimized at $k=\pi$.  The finite grid
samples that point exactly for even $M$; for odd $M$ its nearest points are
$\pi\pm\pi/M$, which proves~\eqref{eq:gaps}.
The curve $q(k)$ is the counterclockwise circle of radius $w$ centred at
$v$.  It encloses zero exactly when $w>v$.  On the critical line it reaches
zero only at $k=\pi$, which belongs to the finite grid exactly for even $M$;
the corresponding $2\times2$ fiber is zero.
\end{proof}

Theorems~\ref{thm:edge} and~\ref{thm:bulk} exhibit a finite separation:
the bulk is topological for $w/v>1$, but the open pair is strictly
hyperbolic only for $w/v>(M+1)/M$.  Equality in the latter condition is a
band-edge taper, not an exponentially localized state.

\begin{theorem}[All hopping faces]\label{thm:faces}
For the open chain the following statements are exact.
\begin{enumerate}
\item If $w=0<v$, there are $M$ dimers and $\pm v$ each has multiplicity $M$.
\item If $v=0<w$, there are $M-1$ dimers, $\pm w$ each has multiplicity
$M-1$, and $A_1,B_M$ span a two-dimensional zero sector.
\item If $v=w=0$, the kernel has dimension $2M$.
\item If $v=w>0$, then
$\spec(H_O)=\{2v\cos(\ell\pi/(2M+1)):1\leq\ell\leq2M\}$, so zero is absent.
\end{enumerate}
For $M=1$, declared separately, the open intercell bond is absent and the
energies are $\pm v$; the periodic wrap merges with the intracell bond and
the energies are $\pm(v+w)$.
\end{theorem}
\begin{proof}
The first three statements follow by listing the isolated sites and bonds.
The fourth is the sine diagonalisation of the uniform path on $2M$ sites;
its denominator $2M+1$ prevents a zero cosine.  The last statement follows
from the two explicit $2\times2$ matrices.
\end{proof}

\begin{theorem}[Entire boundary-safe propagator]\label{thm:propagator}
Let $\sinc z=\sin(z)/z$ with $\sinc0=1$.  Entire matrix functional calculus
gives, on every face in Theorem~\ref{thm:faces},
\begin{equation}\label{eq:propagator}
 e^{-\iu tH_O}=\begin{pmatrix}
 \cos(t\sqrt{TT^*})&-\iu t\,\sinc(t\sqrt{TT^*})T\\
 -\iu t\,\sinc(t\sqrt{T^*T})T^*&\cos(t\sqrt{T^*T})
 \end{pmatrix}.
\end{equation}
It is unitary and satisfies
$\Gamma e^{-\iu tH_O}\Gamma=e^{\iu tH_O}=(e^{-\iu tH_O})^*$.
\end{theorem}
\begin{proof}
The even and odd powers have block forms
\[
 H_O^{2n}=\operatorname{diag}((TT^*)^n,(T^*T)^n),\qquad
 H_O^{2n+1}=\begin{pmatrix}0&(TT^*)^nT\\(T^*T)^nT^*&0\end{pmatrix}.
\]
Separating the exponential series proves~\eqref{eq:propagator}.  Because
cosine and sinc are entire in the squared operators, no inverse or positive
singular-value assumption enters.
Unitarity follows from self-adjointness, and the last identity from chirality.
\end{proof}
\fi

\ifnum\CRevisionRound>1
\section{Quench corollary, evidence, and Route-A boundary}
Consider strictly positive, gapped initial and final hoppings.  Let
$q_i(k),q_f(k)$ use~\eqref{eq:bloch}, and start in the lower initial Bloch
state.

\begin{corollary}[Exact mode zero with finite-grid caveat]\label{cor:quench}
Its final-evolution amplitude is
\begin{equation}\label{eq:loschmidt}
 g_k(t)=\cos(E_f(k)t)+\iu c(k)\sin(E_f(k)t),\qquad
 c(k)=\frac{\operatorname{Re}(\overline{q_i(k)}q_f(k))}
 {|q_i(k)||q_f(k)|}.
\end{equation}
A continuum mode vanishes at a real time if and only if
\begin{equation}\label{eq:cross}
 (v_i-w_i)(v_f-w_f)<0.
\end{equation}
In that case
\begin{equation}\label{eq:criticalmode}
 \cos k_*=-\frac{v_iv_f+w_iw_f}{v_iw_f+w_iv_f},\qquad
 t_n=\frac{\pi/2+n\pi}{E_f(k_*)}.
\end{equation}
A finite ring has such an exact mode zero only if a sampled
$k_m=2\pi m/M$ has this cosine.
\end{corollary}
\begin{proof}
Pauli-matrix exponentiation gives~\eqref{eq:loschmidt}.  A complex zero
requires both $\cos(E_ft)=0$ and $c(k)=0$.  The numerator of $c$ is
\[
 A+B\cos k,\qquad A=v_iv_f+w_iw_f,\quad B=v_iw_f+w_iv_f>0.
\]
It has a root in $[-1,1]$ exactly when $A<B$, since $A,B>0$; but
$A-B=(v_i-w_i)(v_f-w_f)$.  Gapped endpoints make the inequality strict,
giving~\eqref{eq:cross}--\eqref{eq:criticalmode}.  A finite Fourier
decomposition contains only its sampled momenta, proving the last sentence.
\end{proof}

For example, the quench $(v_i,w_i)=(3,1)$ to $(v_f,w_f)=(1,5)$ has
$\cos k_*=-1/2$.  In the audited range, and in general, the finite ring hits
that value exactly when $3$ divides $M$.  Opposite bulk phases alone do not
guarantee an exact zero for an arbitrary finite ring.  This is a
single-particle mode statement, not a many-body dynamical-phase-transition
claim.

\paragraph{Evidence.}
The content-addressed evidence has 55 open-polynomial rows, 33 rational
hyperbolic witnesses, 11 threshold rows, 70 periodic rows containing 595
momentum cells, 33 boundary rows, 30 propagator rows, and six quench rows:
7,161 audited scalar leaves.  A producer-independent checker performs
10,948 checks using exact Sturm sequences and a separate matrix exponential.
SymPy closes 9,181 identities, two isolated runs replay byte for byte, and
53 repaired-hash/parser mutations must fail.  These finite grids are
regressions; Theorems~\ref{thm:edge}--\ref{thm:propagator} and
Corollary~\ref{cor:quench} have the stated analytic domains.

\paragraph{Collision boundary.}
C308 studies a one-site non-Hermitian nonreciprocal Hatano--Nelson skin
chain and excludes topology.  The present object is balanced, Hermitian,
bipartite, and chiral; its distinguishing theorem is the strict finite
bulk--edge threshold separation.  C267 is an infinite uniform-field
Wannier--Stark ladder, C297 a two-site PT dimer, and C138 a metric-graph flux
winding problem.  No resolvent, pseudospectral, or skin-effect claim is
reused here.

\paragraph{Route-A result and nonclaims.}
Under \texttt{NO\_BAD\_EULER\_OR\_ROOT\_NUMBER}, the strict tuple is
\[
 (\mathrm{A0\_FAIL},\mathrm{A1\_FAIL},\mathrm{A2\_FAIL},
  \mathrm{A3\_FAIL},\mathrm{A4\_NATURAL\_QUANTIZATION}).
\]
The cells and Bloch momentum carry no rational-prime payload (A0); the
momentum loop is not a primitive-orbit ledger (A1); the finite determinant
is not an Euler product (A2); and no target functional equation, counting
law, or divisor appears (A3).  The finite Hermitian matrix is already a
natural source quantization (A4), but its lattice levels are not target
zeros and it is not a Hilbert--P\'olya operator.  Route A is rejected and
Route B remains locked.  No target Euler factor, root number, automorphy,
zero correspondence, disorder, interaction, or self-consistent phonon
claim is made.

\paragraph{AI use.}
A generative language model assisted drafting and code scaffolding.  The
displayed derivations, independent recomputation, adversarial tests, and
deterministic artifacts define the audit record.
\fi

\section*{Source lineage}
\begin{thebibliography}{9}
\bibitem{SSH1979}
W. P. Su, J. R. Schrieffer, and A. J. Heeger,
``Solitons in Polyacetylene,'' \emph{Physical Review Letters} 42 (1979),
1698--1701. DOI:
\href{https://doi.org/10.1103/PhysRevLett.42.1698}{10.1103/PhysRevLett.42.1698}.

\bibitem{SSH1980}
W. P. Su, J. R. Schrieffer, and A. J. Heeger,
``Soliton excitations in polyacetylene,'' \emph{Physical Review B} 22
(1980), 2099--2111. DOI:
\href{https://doi.org/10.1103/PhysRevB.22.2099}{10.1103/PhysRevB.22.2099}.

\bibitem{Asboth2016}
J. K. Asb\'oth, L. Oroszl\'any, and A. P\'alyi,
``The Su--Schrieffer--Heeger Model,'' in \emph{A Short Course on
Topological Insulators} (Springer, 2016).
DOI: \href{https://doi.org/10.1007/978-3-319-25607-8_1}%
{10.1007/978-3-319-25607-8\_1}.
\end{thebibliography}
\end{document}
